% Four-page main report; supplementary material begins after the fourth page.
% Replace team/link placeholders before submission. Figures intentionally remain blank.
% Compile from the repository root or this directory with pdflatex, twice.
\documentclass[10pt,letterpaper]{article}
\usepackage[margin=1.9cm]{geometry}
\usepackage{amsmath,booktabs,array,graphicx,hyperref,xcolor}
\hypersetup{hidelinks}
\setlength{\parindent}{0pt}
\setlength{\parskip}{5pt}
\setlength{\textfloatsep}{8pt}
\setlength{\intextsep}{7pt}
\newcommand{\figslot}[2]{\fbox{\parbox[c][#1][c]{0.94\linewidth}{\centering\small\textcolor{gray}{Figure reserved: #2}}}}
\newcommand{\pct}{\%}
\begin{document}
\begin{center}
{\Large\bfseries Two Information Sets, Two Valuation Models}\\[3pt]
{\large SSU Data Dash 2026 --- \textit{TEAM NAME}}\\[2pt]
\textit{FULL NAMES OF ALL TEAM MEMBERS}\\[3pt]
\small Video: \url{https://youtu.be/YOUR-VIDEO} \quad
Code: \url{https://github.com/YOUR-TEAM/YOUR-REPO}
\end{center}

\section*{Summary}
We estimate residential sale prices under two information sets. Off-market OM-V9 learns the price level from property attributes and historical comparable sales; list-aware LA-V22 learns the sale-to-list ratio around the asking price. On 10,069 sales from December 2025 through May 2026, monthly walk-forward validation gives median absolute percentage errors (MdAPE) of \textbf{7.45\%} and \textbf{4.28\%}, respectively. Historical comparables and changes in market pricing matter more than simply substituting one tree algorithm for another. The main weaknesses are expensive, non-standard homes and stale transaction histories; neither model supplies an appraisal or a prediction interval.

\section{The decision, data, and information boundary}
The off-market user is a valuation analyst screening a property before an asking price exists. The list-aware user is assessing a listed property's likely closing price. These are different decisions: the second model has a strong additional price signal, so its lower error is not evidence of a fair algorithm-to-algorithm victory.

The supplied synthetic sales book contains 150,000 transactions closed between January 2021 and May 2026, across six property types, five municipalities, and 145 communities. Condo apartments comprise 51.1\% of the developmental estimation sample. The unlabelled off-market and list-aware books contain 3,974 and 4,064 homes, respectively. No outside data are used. Synthetic coordinates must not be interpreted as real Toronto geography.

Both models use floor area, bedrooms, washrooms, parking, garage, lot measurements, coordinates, property type, style, community, and municipality. Area ranges are represented by their midpoints. Price and floor area are logged; missing lots receive an indicator rather than an invented typical lot. Categories are learned from past data, with unseen categories treated as missing. OM-V9 additionally uses description phrases and a historical text score. LA-V22 uses list price and historical pricing signals.

\textbf{Excluded information.} Off-market never receives list price. Neither model uses days on market, which becomes known after the sale, or the third-party AVM, whose timestamp may postdate the sale and which is absent from the test files. Test sold prices are unavailable and are never used for fitting, selection, or scoring.

\section{Off-market: estimate the price level}
OM-V9 fits LightGBM to log sold price using 35 features. Alongside property attributes, it receives 12-month community/type sold dollars per square foot, city price levels, and earlier sales at the same rounded coordinates and property type. A community cell with fewer than 30 sales falls back to the property type, then the city. Same-pin comparables must close before the home's listing date and exclude its listing identifier.

Description phrases capture condition cues; a TF--IDF ridge model supplies one additional remarks score. Its training scores are constructed in forward blocks so a sale cannot help predict itself. Numbers and community names are masked before text fitting. A separate recent-error correction adjusts the common price level:
\[
\widehat Y_i=\exp\{F_{\rm OM}(x_i)-\delta\},\qquad
\delta=\operatorname{median}_{\text{prior 3 months}}(\log\widehat Y-\log Y).
\]
The calibration predictions come from an earlier fit that excludes those three months. Full correction weight was selected before December 2025. Unlike earlier anchor-gap models, V9 predicts log price directly: historical price levels inform the tree without forcing every prediction to inherit an anchor's error.

\newpage
\section{List-aware: estimate the closing discount}
LA-V22 fits a 23-feature LightGBM to the log sold-to-list ratio:
\[
g_i=\log(Y_i/L_i),\qquad \widehat Y_i=L_i\exp\{F_{\rm LA}(x_i)\}.
\]
This expresses valuation as an adjustment to a strong observed signal. Features describe both the home and its market: peer asking-price levels from historical sales, recent community/type closing discounts, and a listing-month market factor. Community is also a categorical input. A log-linear model and a SARIMAX forecast of its monthly residuals construct the market factor; they are upstream features, not additional dollar predictions averaged with the tree.

The key economic finding is that the sale-to-list gap changes over time. In the fixed-origin development experiment, a log-linear model alone has 7.13\% MdAPE; adding its residual forecast reduces this to 4.78\%. The final tree can also distinguish homes within the same market. A separate citywide haircut was investigated but not retained because its correction did not remain stable across origins.

Both final trees use 700 estimators, 127 leaves, learning rate 0.03, minimum 20 observations per leaf, row and feature subsampling of 0.8, L2 regularization of 1, and seed 0. The list-aware settings were selected using a July--November 2025 inner window and reused for off-market. The aim is to learn nonlinear interactions while limiting repeated hyperparameter search.

\section{Validation: train on the past, score the future}
\textbf{Selection and development.} The inner screen fits through June 2025 and evaluates July--November 2025. A separate fixed-origin developmental fit uses the 139,931 sales closed before December 2025 and evaluates the next 10,069 sales. Its scores are 7.53\% for OM-V9 and 4.43\% for LA-V22. Feature rankings and candidate comparisons on later developmental data also informed development; these six months are therefore not a pristine, untouched final holdout.

\textbf{Headline evaluation.} For each month from December 2025 through May 2026, fit on sales closed strictly before that month and predict its sales. Rebuild the entire pipeline, including lookup tables, categories, text fits, calibration, and time-series features, at each origin. Listing-date lookups additionally restrict comparables to eligible earlier closes. Concatenate the six monthly prediction vectors and calculate metrics over homes; do not average monthly medians. This prevents future closes from entering an earlier month's fit. It does not remove the model-selection uncertainty from repeatedly examining a development window.

We report MdAPE as the primary competition metric, MAPE for average error and tie-breaking, and median signed percentage error for directional bias. With \(e_i=(\widehat Y_i-Y_i)/Y_i\), these are \(100\operatorname{median}|e_i|\), \(100\operatorname{mean}|e_i|\), and \(100\operatorname{median}e_i\). A positive signed value means overvaluation. No confidence interval is asserted for the notebook-derived final-model results; the month path describes temporal stability, not a formal uncertainty bound.

\begin{figure}[h]
\centering
\begin{minipage}{0.48\linewidth}\centering
\figslot{2.35cm}{OM-V9 monthly MdAPE}
\end{minipage}\hfill
\begin{minipage}{0.48\linewidth}\centering
\figslot{2.35cm}{LA-V22 monthly MdAPE}
\end{minipage}
\caption{Monthly walk-forward MdAPE, December 2025--May 2026. Left: OM-V9 versus the 12-month community anchor and V7. Right: LA-V22 versus list price and V21. Each month uses only previously closed sales. The plotted values are documented in Appendix A.}
\label{fig:monthly}
\end{figure}

\newpage
\section{Results: substantial baseline gains, modest final-version gains}
\begin{table}[h]
\centering\small
\begin{tabular}{llrrr}
\toprule
Track & Specification & MdAPE & MAPE & Signed median\\
\midrule
Off-market & Community/type 12-month anchor &15.33\%&20.32\%&+3.19\%\\
& V7 anchor-gap challenger &7.50\%&9.91\%&$-0.50\%$\\
& \textbf{OM-V9 direct log price} &\textbf{7.45\%}&\textbf{9.79\%}&$\mathbf{-0.37\%}$\\
\midrule
List-aware & List price &7.26\%&8.09\%&+5.68\%\\
& V21 without community category &4.36\%&5.46\%&+0.92\%\\
& \textbf{LA-V22 log sold-to-list gap} &\textbf{4.28\%}&\textbf{5.37\%}&\textbf{+0.90\%}\\
\bottomrule
\end{tabular}
\caption{Pooled walk-forward results on the same 10,069 sales. All entries are percentages; lower MdAPE and MAPE are better. The tracks have different information sets. These are validation results, not hidden-test scores.}
\label{tab:results}
\end{table}

OM-V9 reduces MdAPE by 7.88 percentage points, or 51\%, relative to the community anchor. LA-V22 reduces it by 2.98 points, or 41\%, relative to the ask. These are the meaningful baseline improvements. The final-version increments are much smaller: V9 beats V7 by only 0.05 points overall and loses in February; V22 beats V21 by 0.08 points overall, with improvements in all six months. Neither small difference establishes statistical superiority.

MdAPE is a median error, not an accuracy percentage or a guarantee for an individual property. OM-V9's pooled bias is close to zero, but February is 4.44\% low. LA-V22 remains high in every month, with a pooled median overvaluation of 0.90\%. Aggregate scores alone conceal operationally relevant errors.

\section{Where the models fail}
The following segment results come from the \textbf{fixed-origin developmental refits}, not the walk-forward pool. They identify failure patterns and must not be presented as additional walk-forward results.

\begin{figure}[h]
\centering
\begin{minipage}{0.48\linewidth}\centering
\figslot{2.65cm}{OM-V9 error by price decile}
\end{minipage}\hfill
\begin{minipage}{0.48\linewidth}\centering
\figslot{2.65cm}{LA-V22 error by price decile}
\end{minipage}
\caption{Fixed-origin MdAPE by realized sold-price decile, December 2025--May 2026. Left: OM-V9, overall 7.53\%. Right: LA-V22, overall 4.43\%. Deciles are retrospective diagnostics using realized prices, not groups observable before a sale.}
\label{fig:price}
\end{figure}

\textbf{Luxury valuation is the main off-market weakness.} V9's top sold-price decile has 13.58\% MdAPE and $-10.68\%$ median bias, compared with 5.40\% and +2.53\% in the cheapest decile. A common level correction cannot fix opposite biases at the two ends. Unobserved condition, unusual land, and sparse comparable sales are plausible explanations, rather than demonstrated causal findings.

\textbf{Property structure matters more than the broad municipal split.} OM-V9's condo-apartment MdAPE is 6.24\%, versus 9.60\% for detached and 9.84\% for semi-detached homes. Same-pin comparables cover 60\% of developmental validation homes and are more available in condo stock. LA-V22 shows the same broad type pattern: 3.69\% for condo apartments versus 5.52\% for detached and 5.78\% for semi-detached. Co-op results rely on only 48 sales. Appendix B provides the full type table.

\newpage
\section{Recommendations, limits, and responsible use}
\textbf{Use the model that matches the decision.} Apply OM-V9 when no ask exists and LA-V22 when a legitimate listing price is available. Treat each as a point estimate for screening and comparison. For high-value, unusual, or poorly comparable properties, request a manual review of condition and relevant recent sales before making a consequential valuation decision. The retrospective price-decile results motivate that review; they do not supply a ready-made rule using an unknown future sale price.

\textbf{Keep the transaction tape current.} In the V9 May stress test, restricting history to sales closed by January raises MdAPE from 7.39\% to 8.38\% and shifts median bias from +0.98\% to $-5.29\%$. Data freshness is part of the model's operating conditions. Monthly refitting is the evaluated process, but the deployment files extend beyond the May 2026 labelled history. We cannot assume equally fresh comparables or claim that validation error is a guaranteed forecast of hidden-test error.

\textbf{Monitor segments as well as the pool.} Store the model version, prediction, information cutoff, and comparable coverage for every score. As realized sales arrive, track monthly MdAPE, signed bias, and performance by price band and property type. Proposed review triggers from the development documents are monthly MdAPE above 10\% for off-market or 6\% for list-aware, and list-aware median bias outside $\pm2.5\%$. These are review policies, not thresholds validated to optimize a business outcome. Keep the community anchor, list-price benchmark, and immediate predecessor models available as challengers.

\textbf{Be explicit about what remains unproven.} The book is synthetic; results do not establish accuracy on live Toronto transactions. Renovation quality and condition are incompletely observed. No property-specific prediction intervals are provided. Six monthly origins cannot establish long-run robustness or performance in a new market regime. Repeated developmental comparisons create selection risk even when each fitted pipeline respects time order. The hidden test is unlabelled, and its eventual score may be worse. A further untouched time window and a validated uncertainty method would strengthen the evidence.

\section{Reproducibility and submission scope}
The final specifications described here are \textbf{OM-V9 and LA-V22}, documented in the two model-development reports and their source notebooks. The figure scripts reproduce fixed-origin diagnostics; their walk-forward charts use stored notebook outputs and do not independently rerun six monthly fits. Appendix C separates those reproduction paths and lists the exact figure sources.

\textbf{Submission consistency requirement.} At the time of writing, the repository's \texttt{make\_model} entry point selects off-market V6 and list-aware V22; its saved off-market validation result is 8.56\% MdAPE. That is a different model from OM-V9. Before submission, connect the final V9 implementation, rerun monthly validation, and regenerate both CSVs. The V9 scores above are supported by the development notebook, not by the current command-line off-market entry point. This implementation gap must be resolved for the submitted executable and this report to describe the same models.

\textbf{Required delivery.} Replace the team names, video link, and GitHub link on page 1. The repository must contain \texttt{submission\_off\_market.csv} and \texttt{submission\_list\_aware.csv}. After the final implementations are connected, run the standard validation, prediction, and submission checks in Appendix C. The video should explain the information sets, time validation, baseline gains, and one concrete failure mode within five minutes.

\section*{Outside data}
None. All features, diagnostics, and performance statements use the supplied synthetic data. Appendix material documents supporting evidence; the main report ends on this fourth page.

\clearpage
\appendix
\section{Monthly validation and metric provenance}
\small
\begin{table}[h]\centering
\begin{tabular}{lrrrrrr}
\toprule
Month & Homes & OM anchor & OM-V7 & OM-V9 & LA ask & LA-V22\\
\midrule
2025-12 &1238&16.57&7.49&7.42&7.46&3.97\\
2026-01 &1058&15.43&7.44&7.27&7.58&4.16\\
2026-02 &1442&15.61&7.86&8.17&7.15&4.38\\
2026-03 &1816&14.79&7.19&7.10&7.10&4.14\\
2026-04 &2224&15.02&7.64&7.44&7.24&4.37\\
2026-05 &2291&15.28&7.40&7.39&7.25&4.55\\
\midrule
Pooled &10069&15.33&7.50&7.45&7.26&4.28\\
\bottomrule
\end{tabular}
\caption{Walk-forward MdAPE (\%). Pooled statistics are computed over individual predictions, not averages of monthly statistics.}
\end{table}
\begin{table}[h]\centering
\begin{tabular}{lrr}
\toprule
Month & OM-V9 median signed error & LA-V22 median signed error\\
\midrule
2025-12 &+0.63&+1.31\\
2026-01 &$-2.50$&+1.33\\
2026-02 &$-4.44$&+0.20\\
2026-03 &$-1.15$&+0.46\\
2026-04 &+1.65&+1.31\\
2026-05 &+0.98&+0.92\\
\midrule
Pooled &$-0.37$&+0.90\\
\bottomrule
\end{tabular}
\caption{Walk-forward median signed percentage error. A positive value is an overprediction.}
\end{table}

\textbf{Source hierarchy.} Headline and monthly scores are recorded in
\texttt{model-report/OM-V9-off-market-model-report.tex} and
\texttt{model-report/LA-V22-list-aware-model-report.tex}, from the saved walk-forward outputs of \texttt{notebooks/off\_market\_models.ipynb} and \texttt{notebooks/list\_aware\_models.ipynb}. Fixed-origin figures and segment tables come from fresh developmental refits. These have OM-V9 MdAPE/MAPE/bias of 7.53/9.77/$-2.10$\% and LA-V22 values of 4.43/5.53/+1.49\%. The saved command-line validation report currently describes V6/V22 and is not substituted for V9 evidence.

\textbf{Candidate selection caveat.} Several alternatives were inspected on the December--May development window. Walk-forward ensures out-of-time fitting at each origin, but the final candidate was not chosen without reference to that window. We report this distinction to avoid describing the six-month score as a fully independent final test.

\section{Supporting error analysis}
\begin{table}[h]\centering
\begin{tabular}{lrrrrr}
\toprule
Type & Homes & OM MdAPE & OM bias & LA MdAPE & LA bias\\
\midrule
Condo apartment &4843&6.24&$-1.54$&3.69&+1.65\\
Condo townhouse &616&5.98&$-1.42$&4.43&+1.45\\
Co-op apartment &48&10.67&$-1.47$&5.47&+3.60\\
Detached &3123&9.60&$-2.33$&5.52&+1.63\\
Freehold townhouse &405&9.19&$-5.14$&5.58&+0.99\\
Semi-detached &1034&9.84&$-4.34$&5.78&$-0.01$\\
\bottomrule
\end{tabular}
\caption{Fixed-origin developmental errors (\%). These rows are not walk-forward segment estimates. Bias is median signed percentage error.}
\end{table}

\clearpage
\section{Reproduction and figure placement guide}
\normalsize
All requested picture positions intentionally contain blank boxes. The captions are final prose; replace each box with an \texttt{\string\includegraphics} command while preserving its allocated height. All figure paths below are relative to the repository root. From the \texttt{report} directory, prefix each path with \texttt{../}.

\begin{table}[h]\centering\small
\begin{tabular}{p{3cm}p{10.8cm}}
\toprule
Position & Exact image file\\
\midrule
Figure 1, left & \texttt{model-report/figures/om\_walkforward\_mdape.png}\\
Figure 1, right & \texttt{model-report/figures/walkforward\_mdape.png}\\
Figure 2, left & \texttt{model-report/figures/om\_error\_by\_decile.png}\\
Figure 2, right & \texttt{model-report/figures/error\_by\_decile.png}\\
Figure 3, left & \texttt{model-report/figures/om\_error\_by\_type.png}\\
Figure 3, right & \texttt{model-report/figures/error\_by\_type.png}\\
Figure 4, left & \texttt{model-report/figures/om\_calibration.png}\\
Figure 4, right & \texttt{model-report/figures/calibration\_ratio.png}\\
\bottomrule
\end{tabular}
\caption{Placement guide. Figure 1 uses the walk-forward plots; Figures 2--4 use fixed-origin developmental diagnostics.}
\end{table}

Generate the figure files from the repository root:
\begin{quote}\ttfamily\small
py -3 model-report/reproduce\_v9\_figures.py\\
py -3 model-report/reproduce\_v22\_figures.py
\end{quote}
The scripts refit the fixed-origin models and write \texttt{om\_v9\_diagnostics.json} and \texttt{v22\_diagnostics.json}. They do not independently reproduce the monthly refits. To reproduce those results, rerun the walk-forward sections in the two source notebooks, with their frozen final specifications and dates.

After connecting V9 to the command-line interface, the intended submission workflow is:
\begin{quote}\ttfamily\small
pip install -r requirements.txt\\
python run.py validate\\
python run.py report\\
python run.py predict\\
python run.py check
\end{quote}
Compare regenerated metrics and submission-model versions with this report before delivery. The generic generated results table must not silently reintroduce V6 values into a V9 report. The present document uses an explicit final-specification results table for that reason.

\textbf{Layout note.} Explicit page breaks allocate four pages to the main narrative. Appendix begins on physical page 5. Replacing the figure boxes with images of the same height should preserve this allocation. Team details and actual links must still be completed before submission.

\clearpage
\section{Supplementary diagnostic figures}
\begin{figure}[h]\centering
\begin{minipage}{0.48\linewidth}\centering
\figslot{4.0cm}{OM-V9 error by property type}
\end{minipage}\hfill
\begin{minipage}{0.48\linewidth}\centering
\figslot{4.0cm}{LA-V22 error by property type}
\end{minipage}
\caption{Fixed-origin MdAPE by property type. Ground-oriented homes are less accurately valued than condo apartments in both tracks. Co-op estimates are based on only 48 validation sales. Dashed reference lines show each model's fixed-origin overall MdAPE.}
\end{figure}
\begin{figure}[h]\centering
\begin{minipage}{0.48\linewidth}\centering
\figslot{4.0cm}{OM-V9 predicted-price calibration}
\end{minipage}\hfill
\begin{minipage}{0.48\linewidth}\centering
\figslot{4.0cm}{LA-V22 close-ratio calibration}
\end{minipage}
\caption{Fixed-origin calibration diagnostics with different axes and targets. Left: OM-V9 median predicted and realized prices within predicted-price bins; the upper end reveals underpricing. Right: LA-V22 predicted and realized sold-to-list ratios. These plots describe empirical calibration in the development window and do not establish future coverage or a prediction interval.}
\end{figure}

\section{Selected methodological references}
\small
Friedman, J. (2001). Greedy function approximation: a gradient boosting machine. \textit{Annals of Statistics}, 29(5), 1189--1232.

Ke, G. et al. (2017). LightGBM: A highly efficient gradient boosting decision tree. \textit{Advances in Neural Information Processing Systems}, 30.

Rosen, S. (1974). Hedonic prices and implicit markets: Product differentiation in pure competition. \textit{Journal of Political Economy}, 82(1), 34--55.
\end{document}

